Multiplying a row of a mixed-integer program by a positive number preserves it exactly and
changes the numbers a solver sees. We ask what that change buys and what it still guarantees.
On 306 generated hydrothermal-dispatch models, holding each model's LP basis fixed across its
representations, geometric-mean row scaling lowers its condition number on every
model, by six to eight orders of magnitude in the median of each class. An apparent block-metadata effect is a coupling-kernel effect:
geometric-mean centering makes the extrema-based block statistic identically one, so the
block-relative rule becomes Euclidean normalization, which a control without ownership
information reproduces; the idealized coupling problem has optimal factor
$(1+\rho^2)^{-1/2}$. Residual budgets stated in application units, by contrast, survive
consistent re-emission, and we characterize exactly the row factors compatible with a budget
and coefficient limits. In reconstructed dispatch runs with Gurobi, scaling used 30--32\% less
solver work on the 14 SG-Ter-Mer instances every policy completed, while one instance timed out
under every scaled policy and reversed the failure-sensitive ordering, and reruns reproduce the gain with a second seed but not a second solver. In an exactly enumerated stress test, plain centering lost
verified optima on all three solvers, adding the budget to that same rule lost none, and whether scaling helped at all differed across the three
pipelines. A representation can be numerically better without being operationally safer, unless
what safer means is stated in the units of the original model.